\section{The scalar orbital germ from a Hessenberg determinant}
\label{sec:theta-scalar-germ}

We retain $K,R,\psi$ and the Haar normalization of the preceding
sections.  In particular $K$ is a finite product of finite extensions
of $\mathbb Q_p$, $R$ is a compact open subring, and the annihilator
of $R$ is $R$.  An absolute value $|a|$ denotes the modulus of
multiplication by $a$ on the additive group of $K$.  It is the
product of the usual local-field moduli.  All product spaces carry
product Haar measure.  The integer $p$ tends to zero in every
factor, and the lattices $p^aR$, $a\in\mathbb Z$, are cofinal among
neighborhoods of zero.

The purpose of this section is to prove the scalar two-block
orbital germ.  The essential algebraic step is a change from
unipotent coordinates to the first row of an upper Hessenberg
matrix.  Its determinant determines the one remaining diagonal
entry.  Only its expansion through first order is needed; local
constancy makes the resulting formula exact.

\subsection{Test functions and orbital sums}

Fix $m\ge1$, and write $s=m+1$ and $d=m(m+1)/2$.
Let $L,\Lambda\subset K$ be compact open additive subgroups,
and let $Q\subset K^\times$ be compact.  A permitted test function
$\Phi:M_s(K)\to\mathbb C$ satisfies the following conditions:
\begin{enumerate}
\item $\Phi(aZb)=\Phi(Z)$ for integral lower unipotent $a$ and
integral upper unipotent $b$ of size $s$.
\item If $\Phi(Z)\ne0$, then every entry of $Z$ belongs to $L$.
\item If $\det Z\in Q\cup(-Q)$ and $Y\in M_s(\Lambda)$, then
$\Phi(Z+Y)=\Phi(Z)$.
\end{enumerate}
The neighborhoods constructed below depend on $K,R,\psi,m,L,
\Lambda,Q$, and not on the values of $\Phi$.

Let $J_j$ denote the reversing permutation matrix of size $j$.
For a block configuration $w$ of size $s$, define
\[
 \mathcal O_\Phi(w)=
 \sum_{q\in L_w(R)\backslash L_w(K)\times U_s(K)/U_s(R)}
       \Phi(x_qZ(w)y_q)
       \psi\left(\sum_i(x_q)_{i+1,i}
                    +\sum_i(y_q)_{i,i+1}\right).
\]
Here the lower group and representatives have the definitions
used in \eqref{orb-actual-kernel}.  Every sum in this section is
finite.  Indeed the block Gaussian elimination proof of
Lemma~\ref{orb-finite-orbits} applies with the compact set $L$
in place of $R$: bounded product entries and fixed nonzero
diagonal-block determinants bound both unipotent factors.
Invariance of $\Phi$ and triviality of $\psi$ on $R$ prove
independence of representatives.

We abbreviate
\[
 \mathcal O_\Phi(t,c)=\mathcal O_\Phi(\operatorname{diag}(tJ_m,c)),
 \qquad
 \mathcal O_\Phi(zJ_s)=\mathcal O_\Phi(\text{one block }zJ_s).
\]
The integrands below are locally constant on their relevant
determinant sets.  Condition (3), continuity of the matrix maps,
and the compact support just proved show that all Haar integrals
are absolutely convergent.  No global regularity of $\Phi$
outside these sets is needed.

\subsection{The first-row determinant calculation}

Use indices $0,1,\ldots,m$.  For $\rho\in K$, $x=(x_1,\ldots,x_m)$,
and $y=(y_{ij})_{0\le i<j\le m}$, let
$H_t(\rho,x,y)$ be the matrix with diagonal
$(\rho,x_1,\ldots,x_m)$, entries $y_{ij}$ above the diagonal,
entries $t$ immediately below it, and all other entries zero.
Thus $H_t$ is upper Hessenberg.  Let $C_t(x,y)$ be its trailing
$m$ by $m$ principal submatrix; it does not depend on $\rho$ or
on the first-row variables $y_{0j}$.

For $x\in(K^\times)^m$ and $D\in K^\times$, put
\begin{align}
 P(x)&=\prod_{i=1}^m x_i,&
 x_0(D,x)&=D/P(x),\\
 g_D(x,y)&=\frac{y_{01}}{x_1}
     +x_0(D,x)\sum_{i=1}^{m-1}
                       \frac{y_{i,i+1}}{x_ix_{i+1}},&
 f_D(x)&=\sum_{i=1}^m x_i+x_0(D,x).
\end{align}
The empty sum is zero when $m=1$.

\begin{lemma}\label{sg-determinant-expansion}
Fix compact sets of allowed $D\in K^\times$,
$x\in(K^\times)^m$, and $y\in K^d$.  For all sufficiently small
$t$, uniformly over these sets,
\begin{equation}\label{sg-tail-modulus}
 C_t(x,y)\text{ is invertible},\qquad
                  |\det C_t(x,y)|=|P(x)|.
\end{equation}
There is then a unique $\rho=\rho_t(D,x,y)$ with
$\det H_t(\rho,x,y)=D$, and
\begin{equation}\label{sg-rho-expansion}
 \rho_t(D,x,y)=x_0(D,x)+t g_D(x,y)+t^2 E_t(D,x,y),
\end{equation}
where the values of $E_t$ are contained in one compact subset
of $K$.  The neighborhood and this bound are independent of
the individual parameters in the specified compact sets.
\end{lemma}

\begin{proof}
The determinant is linear in $\rho$, with coefficient
$\det C_t$.  At $t=0$ that coefficient is $P(x)$, which is
bounded away from zero in every field factor.  The ratio
$\det C_t/P(x)$ is therefore uniformly in $1+pR$ for small $t$.
Such a ratio is a unit of $R$, with inverse given by a convergent
geometric series, and hence has modulus one.  This proves
\eqref{sg-tail-modulus} and uniqueness of $\rho$.

The coefficient of $t$ in the Hessenberg determinant is
\begin{equation}\label{sg-det-linear}
 -\sum_{i=0}^{m-1}y_{i,i+1}
                 \prod_{j\ne i,i+1}d_j,
 \qquad (d_0,d_1,\ldots,d_m)=(\rho,x_1,\ldots,x_m).
\end{equation}
For completeness, in the permutation expansion a term using
exactly one entry below the diagonal must use an adjacent
transposition.  All other positions must be diagonal: the
sum of column minus row indices is zero, so the sole negative
displacement $-1$ permits exactly one positive displacement
$+1$.  The sign of that transposition is negative, proving
\eqref{sg-det-linear}.  Terms using two or more subdiagonal
entries are divisible by $t^2$.

Put $S=\sum_{i=1}^{m-1}y_{i,i+1}/(x_ix_{i+1})$.
There are polynomial remainders, uniformly bounded on the
compact parameter sets for small $t$, such that
\[
 \det C_t=P(1-tS)+t^2 A_t,
 \qquad
 \det H_t(0,x,y)=-t y_{01}P/x_1+t^2 B_t.
\]
Solving the linear equation for $\rho$ gives
\[
 \rho_t=
 \frac{D+t y_{01}P/x_1-t^2B_t}{P(1-tS)+t^2A_t}.
\]
The denominator divided by $P$ is uniformly near one.
Expanding its inverse by the geometric series, through the
linear term, gives exactly \eqref{sg-rho-expansion}.
All remainders are bounded uniformly because $P^{-1}$, the
polynomial coefficients and the inverse of that near-one
denominator are bounded on the stated compact sets.  This also
proves the assertion when $m=1$, where the higher remainders
in the determinant vanish.
\end{proof}

The next observation ensures that this calculation covers every
contributing orbit, including ones for which the diagonal entries
were not initially assumed invertible.

\begin{lemma}\label{sg-nonzero-diagonals}
Fix a compact $Q_0\subset K^\times$ and a compact open additive
subgroup $L$.  There is a compact open set $E\subset K^\times$
and a neighborhood of zero such that the following holds.
If $t$ belongs to that neighborhood, all entries of
$\rho,x,y$ belong to $L$, and
\[
                    \det H_t(\rho,x,y)\in Q_0,
\]
then every $x_i$ belongs to $E$.
The set $E$ can also be chosen to contain every $x_i$ for which
$x_0(D,x),x_1,\ldots,x_m\in L$ and $D\in Q_0$.
\end{lemma}

\begin{proof}
Choose $a\ge0$ such that $L$, $Q_0$ and $Q_0^{-1}$ are
contained in $p^{-a}R$.  Write the determinant as
\[
             D=\rho\prod_i x_i+t B.
\]
For $t\in R$ all values of the polynomial $B$ lie in one
fixed compact subset of $K$.  Consequently, for small $t$,
$tB/D\in pR$, uniformly.  It follows that
$\rho\prod_i x_i=D(1-tB/D)$ is invertible, and hence so is
every $x_i$.  Moreover
\[
 x_i^{-1}=\frac{\rho\prod_{j\ne i}x_j}{D}
                           (1-tB/D)^{-1}
             \in p^{-(m+1)a}R.
\]
The same bound holds when $\rho=x_0(D,x)$ and $t=0$.
Thus one may take
\[
 E=\{u\in L\cap K^\times:u^{-1}\in p^{-(m+1)a}R\}.
\]
This set is open and compact in $K^\times$: both $u$ and its
inverse are bounded, so no sequence in it can approach a
nonunit in any field component.  It is a finite union of
cosets of some sufficiently small additive lattice.
\end{proof}

\subsection{Construction of the amplitude}

Let $Q_0=Q\cup(-Q)$.  For $D\in Q_0$, define
\begin{equation}\label{sg-amplitude}
 A_D(x)=|P(x)|^{-1}\int_{K^d}
   \Phi\bigl(J_sH_0(x_0(D,x),x,y)\bigr)
                          \psi(g_D(x,y))\,dy
 \quad(x\in(K^\times)^m),
\end{equation}
and set $A_D(x)=0$ otherwise.

\begin{lemma}\label{sg-amplitude-uniform}
There are a compact open $E\subset K^\times$ and an additive
lattice $\Lambda_1\subset K$, depending only on the fixed
data, such that every amplitude \eqref{sg-amplitude} is
supported in $E^m$ and invariant under translations by
$\Lambda_1^m$.  The inner integrands are supported in $L^d$.
All these assertions are uniform in $D\in Q_0$ and in every
permitted $\Phi$.
\end{lemma}

\begin{proof}
If the matrix value of $\Phi$ is nonzero, all its diagonal
and upper entries before reversal belong to $L$.  Their
diagonal product is $D$.  Lemma~\ref{sg-nonzero-diagonals},
at $t=0$, bounds $x$ in $E^m$, and condition (2) bounds $y$
in $L^d$.

Choose an additive lattice small enough to preserve $E$ and
$L$ under translation.  On the compact set
$Q_0\times E^m\times L^d$, the rational functions $P^{-1}$,
$x_0$ and $g_D$ have bounded denominators and are uniformly
continuous.  Shrink the lattice so that translating any of
the $x$ or $y$ coordinates by it has the following effects:
$|P|$ is unchanged; every entry of the change in
$J_sH_0(x_0,x,y)$ lies in $\Lambda$; and the change in
$g_D$ lies in $R$.  These choices are uniform because the
parameter set is compact, and $D$ is held fixed during each
translation.  More explicitly, the differences in reciprocal
coordinates satisfy
\[
 (x_i+u)^{-1}-x_i^{-1}=-u\bigl(x_i(x_i+u)\bigr)^{-1};
\]
their denominators remain in fixed bounded sets.  Products
and sums of these identities give each asserted uniform
continuity bound.

The determinant of the unperturbed matrix is
$(-1)^{s(s-1)/2}D$, which belongs to $Q\cup(-Q)$.
Thus condition (3) makes the matrix factor unchanged.
The character factor is unchanged because $\psi$ is trivial
on $R$.  Outside $E^m\times L^d$, both values vanish because
that set is stable under the selected translations.  This
proves a common translation invariance for the full inner
integrand.  Integrating in $y$ proves the claim for $A_D$.
\end{proof}

\subsection{The exact oscillatory integral}

\begin{theorem}\label{sg-oscillatory-reduction}
For all sufficiently small $t\in K^\times$, uniformly in
$c\in K^\times$ with $ct^m\in Q$ and in permitted $\Phi$,
put $D=(-t)^mc$.  Then
\begin{equation}\label{sg-oscillatory-identity}
 \mathcal O_\Phi(t,c)
   =|t|^{-d}\int_{(K^\times)^m}
                         A_D(x)\psi(f_D(x)/t)\,dx.
\end{equation}
For every $z\in K^\times$ with $z^s=D$,
\begin{equation}\label{sg-critical-amplitude}
 A_D(z,\ldots,z)
       =|z|^{m(m-1)/2}\mathcal O_\Phi(zJ_s).
\end{equation}
\end{theorem}

\begin{proof}
For the two-block matrix $W=\operatorname{diag}(tJ_m,c)$,
the lower group consists of matrices with identity diagonal
blocks and a single arbitrary last row below the diagonal.
Write this row as $(\eta_{m-1},\ldots,\eta_0)$.
Its simple-root phase is $\psi(\eta_0)$.
Write an upper unipotent matrix as $U=I+(u_{ij})_{i<j}$.
Product additive Haar measure in the $u_{ij}$ is Haar measure
on this group, normalized by $\operatorname{vol}(U_s(R))=1$.
Indeed left and right translations are triangular affine
changes in these coordinates with diagonal entries one;
successive one-variable translations preserve product Haar
measure.  The last-row coordinates likewise have measure
normalized by $\operatorname{vol}(R^m)=1$.

The summand is invariant under the corresponding integral
subgroups.  Its support lies in finitely many cosets and each
has volume one.  Hence the orbital sum equals the double
Haar integral of that summand.  Reverse the rows of $LWU$.
The resulting matrix has first row
\[
 t\sum_{i=0}^{m-1}\eta_i U_{i,*}+c e_m^{\,t},
\]
and row $i+1$ equal to $tU_{i,*}$ for $0\le i<m$.
Its determinant is $(-t)^mc=D$: subtract the indicated
linear combination of the last $m$ rows from the first row
and expand along the remaining entry $c$.

Now put $n_{ij}=tu_{ij}$.  This contributes the Jacobian
$|t|^{-d}$.  Set
\[
 x_i=n_{i-1,i}\quad(1\le i\le m),\qquad
 y_{ij}=n_{i-1,j}\quad(1\le i<j\le m).
\]
These are all the scaled upper entries, merely reordered.
The remaining first-row entries are denoted
$\tau_j=y_{0j}$, $1\le j\le m$, and its first diagonal
entry is $\rho=t\eta_0$.  Thus the reversed matrix is
$H_t(\rho,x,y)$ and its phase is
\begin{equation}\label{sg-coordinate-phase}
 \psi\left(\eta_0+\sum_{i=0}^{m-1}u_{i,i+1}\right)
              =\psi\left((\rho+\sum_{i=1}^m x_i)/t\right).
\end{equation}

If the matrix value of $\Phi$ is nonzero, all entries
$\rho,x,y$ belong to $L$.  Its determinant is in $Q_0$.
Lemma~\ref{sg-nonzero-diagonals} therefore restricts $x$ to
a fixed compact open $E^m$, and the upper entries below
the first row lie in $L$.  The same restrictions contain
the support of the amplitude in \eqref{sg-amplitude}.
Outside these restrictions both relevant integrands vanish.

Fix such $x$ and the upper entries below the first row.
The affine map $\eta\mapsto\tau$ has coefficient matrix
\[
             \bigl(tU_{i,j}\bigr)_{1\le j\le m,\ 0\le i<m}
                     =C_t(x,y)^{\,t}.
\]
For small $t$ it is invertible by
Lemma~\ref{sg-determinant-expansion}, and its determinant
has modulus $|P(x)|$.  It is therefore a bijection of all
$K^m$, with $d\eta=|P(x)|^{-1}d\tau$.
The entry $\rho$ after this change is the unique
$\rho_t(D,x,y)$ in that lemma, since the determinant remains
$D$.  In particular, this is a change on the whole first-row
space, not a selection of a branch or a restriction to an
unspecified neighborhood.

When one of the first-row coordinates lies outside $L$,
both matrix test functions vanish.  For all the remaining
$y\in L^d$, the expansion \eqref{sg-rho-expansion} is uniform.
Make $t$ small enough that
\[
 t\in\Lambda,\qquad
 \rho_t-x_0\in\Lambda,\qquad tE_t\in R.
\]
The first two conditions and test-function condition (3),
applied at $J_sH_0(x_0,x,y)$, give
\[
 \Phi(J_sH_t(\rho_t,x,y))
                  =\Phi(J_sH_0(x_0,x,y)).
\]
The third and \eqref{sg-coordinate-phase} give
\[
 \psi\left((\rho_t+\sum_i x_i)/t\right)
                 =\psi(f_D(x)/t)\psi(g_D(x,y)).
\]
Substitute these equalities into the changed integral.
Its two Jacobian factors are $|t|^{-d}$ and $|P(x)|^{-1}$.
Fubini's theorem, applicable to these compactly supported
locally constant functions, gives exactly
\eqref{sg-oscillatory-identity}.

Finally let $x_1=\cdots=x_m=z$ with $z^s=D$.  Then
$x_0=z$ and $g_D(x,y)=z^{-1}\sum_i y_{i,i+1}$.
Substitute $y_{ij}=zv_{ij}$ in \eqref{sg-amplitude}.
The factor $|P|^{-1}$ is $|z|^{-m}$ and the Jacobian is
$|z|^d$.  The matrix becomes $zJ_sU$, where $U$ is upper
unipotent with upper entries $v_{ij}$, and its phase is the
usual $\psi(\sum_i v_{i,i+1})$.  For a single block the lower
group is trivial.  The remaining Haar integral is therefore
$\mathcal O_\Phi(zJ_s)$, by the same volume-one coset argument.
As $d-m=m(m-1)/2$, this proves
\eqref{sg-critical-amplitude}.
\end{proof}

\subsection{The scalar germ}

\begin{corollary}\label{sg-scalar-germ}
For all sufficiently small $t\in K^\times$, uniformly in
$ct^m\in Q$ and in all permitted test functions,
\begin{align}\label{sg-germ-formula}
 \mathcal O_\Phi(t,c)
 ={}&|t|^{-m^2/2}|s|^{-1/2}
   \sum_{z^s=(-t)^mc}\gamma_m(zt)|z|^{m^2/2}
                     \psi(sz/t)\mathcal O_\Phi(zJ_s).
\end{align}
There is no exclusion of residue characteristics dividing
$2s$; all neighborhoods may depend on $m$, the local field
product, the order, the character and the stated compact sets.
\end{corollary}

\begin{proof}
Lemma~\ref{sg-amplitude-uniform} supplies one support and one
translation lattice for every $A_D$.  Apply
Theorem~\ref{th-stationary} with these data and $D\in Q_0$.
Then apply \eqref{sg-oscillatory-identity} and
\eqref{sg-critical-amplitude}.  The exponents combine as
\[
 -d+\frac m2=-\frac{m^2}{2},\qquad
 \frac m2+\frac{m(m-1)}2=\frac{m^2}{2}.
\]
This gives precisely \eqref{sg-germ-formula}, with every
normalization unchanged.  Uniformity follows because both
the oscillatory reduction and stationary phase used
neighborhoods chosen from the fixed data before $\Phi,t,c$.
\end{proof}

This proves the scalar two-block orbital germ and its full
compact-parameter uniformity.  Section~\ref{hr-section} supplies
the test function for the marked-cokernel coefficient with fixed
initial blocks.  The following assembly applies the scalar
formula to that function and then to the actual kernel.


\subsection{Removing the fixed initial blocks}
\label{hr-section}

We now prove the reduction of a scalar orbital coefficient to a
function of the last blocks. This is the reduction required before
the stationary-phase calculation. All sums below are finite on their
effective support. The proof retains the specified vector and
inverse-coordinate functional of each cokernel.

Fix a configuration $h$ of total size $a$, put $H=Z(h)$ and
$d=\det H\in K^\times$, and fix a finite marked module
$T=(C,e,\lambda)$. For an integral matrix $B$ of size $b$, let
$T_B$ denote $R^b/BR^b$ with its usual last-vector and
last-inverse-coordinate specifications when $B$ is invertible.
Define
\[
 F_T(B)=
 \begin{cases}
 \#\operatorname{Iso}(T_B,T),&B\in M_b(R)\cap\mathrm{GL}_b(K),\\
 0,&\text{otherwise}.
 \end{cases}
\]
In size zero this uses the zero module and zero specifications.
For an integral singular matrix its cokernel cannot be a finite
module: extension from $R$ to $K$ leaves a nonzero cokernel in at
least one field factor, whereas extension of a finite module is zero:
its cardinality annihilates it and is a nonzero integer invertible in $K$.
Here $K=R[1/p]$, since $R$ is open. Thus assigning zero in the
singular case agrees with counting isomorphisms to the finite target.

Fix $s\ge1$. For $A\in\mathrm{GL}_a(K)$ and rectangular matrices
$X\in M_{s,a}(K)$, $Y\in M_{a,s}(K)$, put
\begin{equation}\label{hr-block-matrix}
 B(A,X,Y,G)=
 \begin{pmatrix}A&AY\\ XA&XAY+G\end{pmatrix}
 =\begin{pmatrix}I&0\\X&I\end{pmatrix}
   \begin{pmatrix}A&0\\0&G\end{pmatrix}
   \begin{pmatrix}I&Y\\0&I\end{pmatrix}.
\end{equation}
Write $\ell(X)=X_{1,a}$ and $r(Y)=Y_{a,1}$ when $a>0$,
and put both expressions equal to zero when $a=0$.

For each integral orbital coset $q\in\mathcal I_h$, choose its
representative matrix $A_q=x_qHy_q$ and phase $a_q$. Define
\[
 \mathcal X_q=
  \{X\bmod M_{s,a}(R):XA_q\in M_{s,a}(R)\},\qquad
 \mathcal Y_q=
  \{Y\bmod M_{a,s}(R):A_qY\in M_{a,s}(R)\}.
\]
These are finite sets. For example, the first is the quotient
$M_{s,a}(R)A_q^{-1}/M_{s,a}(R)$ of two compact open additive
lattices; the denominator is contained in the numerator because
$A_q$ is integral. The second is treated in the same way. Put
\begin{equation}\label{hr-Phi}
 \Phi_{h,T,s}(G)=
 \sum_{q\in\mathcal I_h}a_q
    \sum_{X\in\mathcal X_q}\sum_{Y\in\mathcal Y_q}
        \psi(\ell(X)+r(Y))F_T(B(A_q,X,Y,G)).
\end{equation}
The sets indexing this sum do not depend on $G$ or on $T$.

\begin{lemma}\label{hr-well-defined}
Formula \eqref{hr-Phi} is independent of all the indicated
representatives. If $v\in U_s^-(R)$ and $u\in U_s^+(R)$ are
lower and upper unipotent matrices, respectively, then
\begin{equation}\label{hr-invariance}
                  \Phi_{h,T,s}(vGu)=\Phi_{h,T,s}(G).
\end{equation}
Moreover,
\begin{equation}\label{hr-support}
 \Phi_{h,T,s}(G)\ne0
       \quad\Longrightarrow\quad G\in d^{-1}M_s(R).
\end{equation}
\end{lemma}

\begin{proof}
For any integral lower unipotent matrix $v_0$ and upper
unipotent matrix $u_0$, the replacement $B\mapsto v_0Bu_0$
preserves integrality and invertibility. It preserves the marked
module: the induced cokernel map sends $[x]$ to $[v_0x]$,
the last vector is fixed by $v_0$, and the last row of $u_0^{-1}$
is the last coordinate row. Hence $F_T(v_0Bu_0)=F_T(B)$.

Replacing $X,Y$ by $X+X_0,Y+Y_0$ with integral $X_0,Y_0$
multiplies \eqref{hr-block-matrix} on the left and right by
integral unipotent matrices. The corner phase changes by an
element of $R$, on which $\psi$ is trivial. This proves
independence of the rectangular representatives.
If the representative of $q$ changes $A_q$ to $v_0A_qu_0$,
the bijection $X\mapsto Xv_0^{-1}$, $Y\mapsto u_0^{-1}Y$
identifies the summands. The block matrix is then multiplied
by $\operatorname{diag}(v_0,I_s)$ and
$\operatorname{diag}(u_0,I_s)$. The corner entries and the
phase $a_q$ are unchanged modulo $R$.

For \eqref{hr-invariance}, use the bijections
$X\mapsto vX$ and $Y\mapsto Yu$. Then
\[
 B(A,vX,Yu,vGu)=\operatorname{diag}(I_a,v)
                  B(A,X,Y,G)\operatorname{diag}(I_a,u).
\]
The first row of a lower unipotent matrix and the first column
of an upper unipotent matrix are coordinate vectors. Therefore
$\ell(vX)=\ell(X)$ and $r(Yu)=r(Y)$, which proves the identity.

If the finite sum is nonzero, some block matrix
$B=B(A_q,X,Y,G)$ is integral. Its upper-left block $A=A_q$,
both $XA$ and $AY$, and its lower-right block $XAY+G$ are
integral. Also $\det A=d\in R$ and $\operatorname{adj}(A)$
is integral. The identity
\[
 dG=d(XAY+G)-(XA)\operatorname{adj}(A)(AY)
\]
shows that $dG$ is integral. For $a=0$, $d=1$ and the same
assertion follows directly from $B=G$.
\end{proof}

We next give the perturbation estimate that preserves the
functional as well as the underlying module.

\begin{lemma}\label{hr-perturbation}
Let $B\in\mathrm{GL}_b(K)$, write $\delta=\det B$, and
suppose that $P\in M_b(R)$ and
$P\in\delta^2pM_b(R)$. Then
\[
                         F_T(B+P)=F_T(B)
\]
for every fixed finite marked module $T$.
If $B$ is integral, the identification is the identity on the
ambient module $R^b$: the two image lattices are equal and
$(B+P)^{-1}-B^{-1}$ is integral.
\end{lemma}

\begin{proof}
If $B$ is not integral, adding the integral matrix $P$ cannot
make it integral, so both sides are zero. Assume $B$ integral
and write $P=\delta^2pP_0$ with $P_0$ integral. Put
\[
 A=I+B^{-1}P=I+p\delta\operatorname{adj}(B)P_0.
\]
Since $\delta\in R$, this is a matrix in $I+pM_b(R)$.
Every such matrix is invertible over $R$: its geometric-series
inverse converges, because the powers of a matrix in $pM_b(R)$
have entries in $p^jR$ and these sets tend to zero. Compactness
of $R$ makes it complete. This argument also applies at $p=2$
and to a nonmaximal order.

Now $B+P=BA$, so $(B+P)R^b=BR^b$. Moreover
\[
 (B+P)^{-1}-B^{-1}
   =-pA^{-1}\operatorname{adj}(B)P_0\operatorname{adj}(B)
                                                     \in M_b(R).
\]
The identity on the common cokernel preserves the specified
vector and every inverse-coordinate functional modulo $R$.
It consequently identifies the two sets of marked
isomorphisms to $T$.
\end{proof}

\begin{corollary}\label{hr-local-constancy}
Let $Q\subset K^\times$ be compact. There is an additive
lattice $L\subset R$, depending only on $R,d,Q$, such that
\[
 \Phi_{h,T,s}(G+E)=\Phi_{h,T,s}(G)
 \quad\text{if }\det G\in Q\cup(-Q),\quad E\in M_s(L).
\]
The choice is independent of $T$ and $s$.
\end{corollary}

\begin{proof}
Choose $k\ge0$ such that $p^kd^{-2}q^{-2}\in R$ for every
$q\in Q$, using compactness of the inverse set; take
$L=p^{k+1}R$. In every term of \eqref{hr-Phi}, the full block
matrix has determinant $d\det G$. Replacing $G$ by $G+E$
adds only the lower-right block $E$, an integral perturbation
whose entries divided by $(d\det G)^2$ belong to $pR$.
Lemma~\ref{hr-perturbation} therefore applies term by term.
The indexing sets and phases do not change.
\end{proof}

\begin{proposition}\label{hr-sum-identity}
Let $g$ be any configuration of total size $s$, and let $h\,g$
denote concatenation of their lists of blocks. Let
\[
 \mathcal Q_g=L_g(R)\backslash L_g(K)\times U_s(K)/U_s(R)
\]
be its full orbital coset space. For $q\in\mathcal Q_g$, write
$G_q=x_qZ(g)y_q$ and let $a_q$ be its phase. Then
\begin{equation}\label{hr-orbital-sum}
 \mathcal B(h\,g;T)
           =\sum_{q\in\mathcal Q_g}a_q\Phi_{h,T,s}(G_q).
\end{equation}
The support of the sum on the right is finite. The matrices
$G_q$ in it are not required to be integral.
\end{proposition}

\begin{proof}
Split the coordinates into the first $a$ and last $s$.
Every lower matrix for the concatenated configuration has
a unique expression
\[
 \begin{pmatrix}v_h&0\\Xv_h&v_g\end{pmatrix},
 \qquad v_h\in L_h(K),\quad v_g\in L_g(K),
\]
and every upper matrix has a unique expression
\[
 \begin{pmatrix}u_h&u_hY\\0&u_g\end{pmatrix},
 \qquad u_h\in U_a(K),\quad u_g\in U_s(K).
\]
Choose representatives for the four diagonal cosets. Left
integral multiplication on the lower side then leaves exactly
$X\bmod M_{s,a}(R)$; right integral multiplication on the
upper side leaves exactly $Y\bmod M_{a,s}(R)$.
For clarity, once the diagonal representatives agree, the
quotient of two lower matrices is
$\left(\begin{smallmatrix}I&0\\X'-X&I\end{smallmatrix}\right)$;
the corresponding quotient on the upper side is
$\left(\begin{smallmatrix}I&Y'-Y\\0&I\end{smallmatrix}\right)$.
This proves both existence and uniqueness of the resulting
coset parametrization.

Their orbital matrix is exactly
$B(v_hHu_h,X,Y,v_gZ(g)u_g)$.
The adjacent-entry phase factors into the phases of the first
and last configurations, multiplied by
$\psi(\ell(X)+r(Y))$. Indeed the only adjacent entry crossing
the block boundary is the corner entry, and
$\ell(Xv_h)=\ell(X)$, $r(u_hY)=r(Y)$ by triangularity.
No determinant or measure factor enters this bijection of
discrete cosets.

If the full orbital matrix is integral, its first block is
integral, and $XA$ and $AY$ are integral. Thus the nonzero
terms are exactly among those defining \eqref{hr-Phi}.
There are only finitely many full integral orbital cosets by
Lemma~\ref{orb-finite-orbits}. The bijection just proved
therefore has finite effective support, so it can be regrouped
without an infinite-sum interchange. Grouping first by the
coset of the last configuration gives
\eqref{hr-orbital-sum}.
\end{proof}

Applying Proposition~\ref{hr-sum-identity} with
$g=[m,t][1,c]$ and with $g=[m+1,z]$ gives the two scalar
identities required to remove the fixed initial blocks. The
same function $\Phi_{h,T,m+1}$ occurs in both. It is invariant
under integral lower and upper unipotent multiplication, is
supported in $d^{-1}M_{m+1}(R)$, and is locally constant on
the determinant range $Q\cup(-Q)$ with a lattice independent
of the target module. These are precisely the required
hypotheses on the matrix function in the subsequent
oscillatory reduction.

\begin{remark}[Nonintegral matrices in the last blocks]
Their inclusion in \eqref{hr-orbital-sum} is essential.
Take $K=\mathbb Q_3$, $R=\mathbb Z_3$, the standard character
with $\psi(1/3)=\zeta_3$, one initial block $H=(3)$,
one last block $G=(-1/3)$, and the zero marked target module.
The full orbital matrices are
\[
 \begin{pmatrix}3&3y\\3x&3xy-1/3\end{pmatrix}.
\]
Integrality forces $x=u/3$, $y=v/3$ modulo $R$, with
$uv\equiv1\pmod3$. Their determinant is $-1$, so every
integral matrix has zero cokernel and contributes its phase.
The only pairs are $(u,v)=(1,1),(2,2)$; their sum is
\[
             \zeta_3^2+\zeta_3^4=-1.
\]
The last scalar $-1/3$ is not integral. Restricting the right
side of \eqref{hr-orbital-sum} to integral last-block matrices
would incorrectly give zero. The actual source retains the
full coset sum.
\end{remark}

\begin{remark}[Preserving the module does not preserve the functional]
Take $R=\mathbb Z_p$, $B=(p^2)$ and $P=(p^3)$.
The image ideals of $B$ and $B+P=p^2(1+p)$ are equal, but
their inverse entries differ by
$-1/(p(1+p))\notin R$. The two cyclic cokernels with marked
vector $[1]$ are the same underlying module; any isomorphism
preserving this generator is the identity. Their specified
functionals differ on that generator, so they are not
isomorphic as marked modules. This explains why a perturbation
of size merely $p\det(B)R$ is insufficient. The squared
determinant condition in Lemma~\ref{hr-perturbation} retains
the information that is actually used in the germ formula.
\end{remark}


\subsection{The germ formula for the actual marked-cokernel kernel}
\label{ga-section}

We now combine the initial-block reduction and the scalar
oscillatory calculation.  All objects in the statement are
the actual finite orbital sums defined above.

\begin{theorem}\label{ga-actual-germ}
Fix a block configuration $h$, an integer $m\ge1$, and a
compact set $Q\subset K^\times$.  Put $s=m+1$.
There is a neighborhood $V$ of zero such that for all
$t,c\in K^\times$ with $t\in V$ and $ct^m\in Q$, and every
comparison configuration $w'$, one has
\begin{align}\label{ga-actual-kernel-formula}
 \mathcal K(W_h(t,c),w')
 ={}&|t|^{-m^2/2}|s|^{-1/2}
       \sum_{z^s=(-t)^mc}\gamma_m(zt)|z|^{m^2/2}
          \psi(sz/t)\mathcal K(h[s,z],w').
\end{align}
The size, block decomposition and parameters of $w'$ are
unrestricted.  The neighborhood is chosen before $t,c,w'$.
For a fixed finite collection of possible $h$, one neighborhood
works for every member of that collection.
\end{theorem}

\begin{proof}
Write $d_h=\det Z(h)$.  For each fixed finite marked target
module $T$, Section~\ref{hr-section} constructs the test function
$\Phi_{h,T,s}$.  Lemma~\ref{hr-well-defined} makes it invariant
under integral lower and upper unipotent multiplication and
supports it on $d_h^{-1}M_s(R)$.  Corollary~\ref{hr-local-constancy}
provides one additive lattice on which it is locally constant
when the determinant belongs to $Q\cup(-Q)$.
These bounds are independent of the target module.
They are precisely the three test-function assumptions of
Section~\ref{sec:theta-scalar-germ}.

Apply Corollary~\ref{sg-scalar-germ}.  Proposition~\ref{hr-sum-identity}
identifies its two orbital coefficients, giving
\begin{align}\label{ga-marked-formula}
 \mathcal B(W_h(t,c);T)
 ={}&|t|^{-m^2/2}|s|^{-1/2}
       \sum_{z^s=(-t)^mc}\gamma_m(zt)|z|^{m^2/2}
               \psi(sz/t)\mathcal B(h[s,z];T).
\end{align}
All sums here include every orbital coset on their finite
effective supports, including nonintegral last-block matrices
in the initial-block reduction.

Proposition~\ref{gu-uniform-germ} already shows that these
identities against fixed marked targets imply the germ
uniformly against all comparison configurations.  In the
present proof one can also see the uniformity directly:
the test-function bounds and the scalar-germ neighborhood
were independent of $T$.  For any $w'$, apply
\eqref{ga-marked-formula} to each marked module $C_{Z_{q'}}$
in its finite integral orbital expansion, multiply by
$\overline{a_{q'}}$, and sum.  Formula
\eqref{gu-kernel-coefficients} gives
\eqref{ga-actual-kernel-formula}.  The root sum is finite,
so the interchange introduces no convergence requirement.
This argument does not require the size of $w'$ to equal
the size of $W_h(t,c)$.

Finally, intersect the neighborhoods for finitely many $h$.
This retains a neighborhood of zero and proves the last assertion.
\end{proof}

\begin{corollary}\label{ga-actual-truncation}
For the actual marked-cokernel kernel, the truncated orbital
identity \eqref{orb-exact-cut} holds under the local hypotheses
stated in Theorem~\ref{orb-cut-identity}, without an additional
germ assumption.
\end{corollary}

\begin{proof}
The Gram formula gives Hermitian symmetry.
Proposition~\ref{orb-local-constancy-proved} supplies the
required local constancy, Corollary~\ref{orb-support-proved}
supplies the inverse-parameter support, and
Theorem~\ref{ga-actual-germ} supplies the last hypothesis.
Apply Theorem~\ref{orb-cut-identity}.
\end{proof}

This closes the local orbital germ and truncated orbital identity.
It does not assert a global theta energy bound: the global
arithmetic and automorphic estimates remain distinct statements.


